기룡이는 여러 종류의 진공관을 생산하고 판매하는 공장의 사장이다. 
진공관은 여러 쓰임새가 있지만, 기룡이는 음악 감상을 위한 진공관 앰프에 사용할 진공관을 전문적으로 다루고 있다. 
진공관 앰프는 제품에 따라 모노 출력으로 한 개의 진공관만을 쓰기도 하고, 스테레오 출력을 위해 동일한 종류의 진공관 두 개를 사용하기도 한다. 
따라서 기룡이는 생산된 진공관들을 한 개씩 개별 포장하거나 두 개씩 짝지어 포장해야 한다.

하지만 같은 종류의 진공관이라도 생산할 때마다 미세한 오차가 생길 수 있고, 이 오차가 누적되면 소리의 밸런스가 틀어질 수 있다. 
따라서 최대한 밸런스가 잘 맞도록 하기 위해서는 연속해서 생산된 같은 종류의 진공관 두 개만을 한 쌍으로 포장해야 한다. 

열심히 진공관을 짝지어 포장 중인 기룡이를 구경하던 솔리는 진공관을 낱개로 포장하거나, 같은 종류의 연속한 두 진공관을 한 쌍으로 포장하는 경우의 수가 몇 개인지 궁금하였다.

생산된 $N$개의 진공관의 종류가 주어질 때, 기룡이가 진공관들을 포장할 수 있는 경우의 수를 구해보자. 모든 진공관은 정확히 하나의 포장에 포함되어야 한다.

예를 들어 \t{Mul}이란 진공관이 연속적으로 세 개 생산되었다면, 아래와 같이 세 가지 방법으로 포장할 수 있다.
\begin{enumerate}
  \item \t{(Mul), (Mul), (Mul)}
  \item \t{(Mul, Mul), (Mul)}
  \item \t{(Mul), (Mul, Mul)}
\end{enumerate}